%%% BEGIN fig:dy-two-gluon-exchange
\begin{figure}[htbp]
\centering
\begin{tikzpicture}[fd,
  fc lab/.style={font=\small},
  fc note/.style={font=\footnotesize},
  fc tiny/.style={font=\scriptsize}]
  % =====================================================================
  % (a) box: the two gluons are attached in the same order on both lines
  % =====================================================================
  \begin{scope}
    \draw[gluon] (-0.8,-0.85) -- (-0.8,0.85);
    \draw[gluon] (0.8,-0.85) -- (0.8,0.85);
    \draw[wilson=0.14] (-2.1,0.85) -- (2.1,0.85);
    \draw[wilson=0.14] (-2.1,-0.85) -- (2.1,-0.85);
    \foreach \p in {(-0.8,0.85),(0.8,0.85),(-0.8,-0.85),(0.8,-0.85)} \node[vertex] at \p {};
    \path[mom={$k$},mom sep=0.27cm] (-0.8,-0.85) -- (-0.8,0.85);
    \path[mom'={$K-k$},mom sep=0.27cm] (0.8,-0.85) -- (0.8,0.85);
    \node[fc lab,anchor=east] at (-2.15,0.85) {$v_A,\,\mathbf T_A$};
    \node[fc lab,anchor=east] at (-2.15,-0.85) {$v_B,\,\mathbf T_B$};
    \node[fc tiny] at (0,1.33) {$\dfrac{\ii}{x_A+\ii0}$};
    \node[fc tiny] at (0,-1.35) {$\dfrac{\ii}{-x_B+\ii0}$};
    \node[fc lab] at (0,-2.25) {$J_{\rm box}=\dfrac{2\pi}{ab}\,(\pi+\ii Y)$};
    \node[fc lab] at (0,-3.15) {$c_{\rm box}=X^2$};
    \node[fc lab] at (0,-3.95) {(a)\ 非交叉图（box）};
  \end{scope}
  \node[fc lab] at (3.35,0) {$+$};
  % =====================================================================
  % (b) crossed box
  % =====================================================================
  \begin{scope}[xshift=6.7cm]
    \draw[gluon] (-0.8,-0.85) -- (0.8,0.85);
    % the gluon k passes over the gluon K-k
    \draw[white,line width=7.5pt] (0.2,-0.21) -- (-0.2,0.21);
    \draw[gluon] (0.8,-0.85) -- (-0.8,0.85);
    \draw[wilson=0.14] (-2.1,0.85) -- (2.1,0.85);
    \draw[wilson=0.14] (-2.1,-0.85) -- (2.1,-0.85);
    \foreach \p in {(-0.8,0.85),(0.8,0.85),(-0.8,-0.85),(0.8,-0.85)} \node[vertex] at \p {};
    \path[mom={$k$},mom pos=0.78,mom sep=0.27cm] (0.8,-0.85) -- (-0.8,0.85);
    \path[mom'={$K-k$},mom pos=0.78,mom sep=0.27cm] (-0.8,-0.85) -- (0.8,0.85);
    \node[fc lab,anchor=east] at (-2.15,0.85) {$v_A,\,\mathbf T_A$};
    \node[fc lab,anchor=east] at (-2.15,-0.85) {$v_B,\,\mathbf T_B$};
    \node[fc tiny] at (0,1.33) {$\dfrac{\ii}{x_A+\ii0}$};
    \node[fc tiny] at (0,-1.35) {$\dfrac{\ii}{x_B+\ii0}$};
    \node[fc lab] at (0,-2.25) {$J_{\rm cross}=-\dfrac{2\pi\ii Y}{ab}$};
    \node[fc lab] at (0,-3.15) {$c_{\rm cross}=X^2+\dfrac{C_A}{2}\,X$};
    \node[fc lab] at (0,-3.95) {(b)\ 交叉图（crossed box）};
  \end{scope}
  % =====================================================================
  % sum
  % =====================================================================
  \node[fc lab] at (3.35,-5.15) {$J_{\rm box}\,c_{\rm box}+J_{\rm cross}\,c_{\rm cross}
     =\dfrac{2\pi}{ab}\Big[\underbrace{\pi X^2}_{\mbox{\footnotesize 指数化项}}\;\underbrace{-\,\ii Y\,\tfrac{C_A}{2}\,X}_{\mbox{\footnotesize Regge 相位}}\Big]$};
\end{tikzpicture}
\caption{两条程函线之间的双胶子交换（\cref{prop:dy-eikonal-twogluon}）。两条无限长类时 Wilson 线的速度为$v_A$、$v_B$（快度差$Y>0$，$v_A\cdot v_B=\cosh Y$），色生成元为$\mathbf T_A$、$\mathbf T_B$，时间自左向右；两个胶子取质量$m$（红外正规子）并保留完整传播子，横动量$\mathbf k_T$、$(\mathbf K-\mathbf k)_T$固定，$a=\mathbf k_T^2+m^2$，$b=(\mathbf K-\mathbf k)_T^2+m^2$。
线段旁为两顶点之间的程函传播子（$x_{A,B}=v_{A,B}\cdot k$），其下为纵向积分$J$与色因子$c$（$X=\mathbf T_A\cdot\mathbf T_B$）。
全部$s$道权重在非交叉图中，交叉图只贡献单圈 Regge 相位$\propto C_AY$：横向积分后第一项是指数化项$\frac12(-\ii\chi)^2X^2$，第二项是$\ii\times$实数$\times$厄米色矩阵；因而$p_{\rm opp}=-\ii Y/\pi$，${\rm Re}\,p_{\rm opp}=0$，$\order{g^4}$的幺正性没有残差（VF-REF-9，VF-REF-10）。}
\label{fig:dy-two-gluon-exchange}
\end{figure}
%%% END fig:dy-two-gluon-exchange
